\section{Dataset Construction: The \textsc{RippleEval} Benchmark}
\label{sec:dataset}

To strictly analyze how noise propagates through LLM internal knowledge, we construct \textbf{\textsc{RippleEval}}, a controlled knowledge graph benchmark consisting of paired factual and counterfactual triplets. Unlike existing benchmarks that mix 1-to-N relationships, \textsc{RippleEval} enforces strict \textbf{functional (1-to-1) constraints} (e.g., \textit{CapitalOf}, \textit{DateOfBirth}) to eliminate ambiguity as a confounding factor in error propagation analysis.

\paragraph{Graph Expansion and Verification.}
We employ a ``Generate-then-Verify'' pipeline to construct the factual ground truth graph, denoted as $\mathcal{G}_{fact}$.
\begin{enumerate}
    \item \textbf{Seed Initialization:} We anchor the graph with high-confidence seed triplets verified against Wikidata (e.g., popular entities from Wikipedia).
    \item \textbf{Constraint-Aware Expansion:} Using \texttt{gpt-4-turbo}, we expand the graph by predicting 1-to-1 relationships. A \textbf{Strict Constraint Filter} discards any relations violating the functional property (e.g., rejecting \textit{AuthorOf} for authors with multiple works) to ensure a deterministic propagation path.
    \item \textbf{Fact Verification:} All candidate triplets are verified against Wikidata using exact string matching and alias lookup. The final $\mathcal{G}_{fact}$ contains \textbf{10,379 nodes} and \textbf{16,948 edges}.
\end{enumerate}

\paragraph{Counterfactual Injection (The Noise Set).}
To simulate the ``noise'' described in our problem definition, for each target triplet $(s, r, o) \in \mathcal{G}_{fact}$, we generate a plausible but incorrect counterfactual object $o'$ (e.g., changing \textit{Paris} to \textit{Lyon} for France's capital). These form the perturbation set $\mathcal{G}_{noise}$, which is used to inject specific errors into the model during the updating phase.

\paragraph{Topological Stratification.}
We stratify the nodes in $\mathcal{G}_{fact}$ based on \textbf{In-Degree Popularity} to analyze structure-dependent behaviors:
\begin{itemize}
    \setlength\itemsep{0em} % Save space for ACL Short
    \item \textbf{High-Popularity (Hubs):} The top 5\% of nodes by in-degree, representing widely connected entities.
    \item \textbf{Low-Popularity (Tails):} The bottom 50\% of nodes, representing long-tail knowledge.
\end{itemize}
This stratification allows us to quantitatively compare error propagation rates between ``Hub'' and ``Tail'' regions, facilitating the analysis of the \textit{Hub-Failure} and \textit{Tail-Confidence} phenomena.
